\uuid{94JA}
\titre{Implication/équivalence}
\auteur{bodin ; hanani ; mzari}
\organisation{exo7}
\licence{CC BY-NC-SA 4.0}
\datecreate{2020-01-20}
\niveau{L1}
\module{Algèbre}
\chapitre{Logique, ensemble, raisonnement}
\sousChapitre{Logique}
\difficulte{2}
\theme{Logique, implication, équivalence}
\exo7qcmid{lille-1.logique.009}

\contenu{
\begin{question}
\qclassification{100.01}
\qtags{motcle=implication/équivalence}
\qkeeporder
Par quoi peut-on compléter les pointillés pour avoir une assertion vraie ?
$$|x^2| < 5 \quad \ldots \quad -\sqrt{5} < x < \sqrt{5}$$
\begin{answers}
    \good{$\Longleftarrow$}

    \good{$\implies$}

    \good{$\iff$}

    \bad{Aucune des réponses ci-dessus ne convient.}
\end{answers}
\begin{explanations}
C'est une équivalence, donc en particulier les implications dans les deux sens sont vraies !
\end{explanations}
\end{question}
}
